%%% Hlavní soubor. Zde se definují základní parametry a odkazuje se na ostatní části. %%%

%% Verze pro jednostranný tisk:
% Okraje: levý 40mm, pravý 25mm, horní a dolní 25mm
% (ale pozor, LaTeX si sám přidává 1in)
\documentclass[12pt,a4paper]{report}
\setlength\textwidth{145mm}
\setlength\textheight{247mm}
\setlength\oddsidemargin{15mm}
\setlength\evensidemargin{15mm}
\setlength\topmargin{0mm}
\setlength\headsep{0mm}
\setlength\headheight{0mm}
% \openright zařídí, aby následující text začínal na pravé straně knihy
\let\openright=\clearpage

%% Pokud tiskneme oboustranně:
% \documentclass[12pt,a4paper,twoside,openright]{report}
% \setlength\textwidth{145mm}
% \setlength\textheight{247mm}
% \setlength\oddsidemargin{15mm}
% \setlength\evensidemargin{0mm}
% \setlength\topmargin{0mm}
% \setlength\headsep{0mm}
% \setlength\headheight{0mm}
% \let\openright=\cleardoublepage

%% Pokud používáte csLaTeX (doporučeno):
\usepackage{czech}
%% Pokud nikoliv:
%\usepackage[czech]{babel}
%\usepackage[T1]{fontenc}

%% Použité kódování znaků: obvykle latin2, cp1250 nebo utf8:
\usepackage[utf8]{inputenc}

%% Ostatní balíčky
\usepackage{graphicx}
\usepackage{amsthm}
\usepackage{listings}

%% Balíček hyperref, kterým jdou vyrábět klikací odkazy v PDF,
%% ale hlavně ho používáme k uložení metadat do PDF (včetně obsahu).
%% POZOR, nezapomeňte vyplnit jméno práce a autora.
\usepackage[ps2pdf,unicode]{hyperref}   % Musí být za všemi ostatními balíčky
\hypersetup{pdftitle=Název práce}
\hypersetup{pdfauthor=Bc. Ondřej Odcházel}

%%% Drobné úpravy stylu

% Tato makra přesvědčují mírně ošklivým trikem LaTeX, aby hlavičky kapitol
% sázel příčetněji a nevynechával nad nimi spoustu místa. Směle ignorujte.
\makeatletter
\def\@makechapterhead#1{
  {\parindent \z@ \raggedright \normalfont
   \Huge\bfseries \thechapter. #1
   \par\nobreak
   \vskip 20\p@
}}
\def\@makeschapterhead#1{
  {\parindent \z@ \raggedright \normalfont
   \Huge\bfseries #1
   \par\nobreak
   \vskip 20\p@
}}
\makeatother

% Toto makro definuje kapitolu, která není očíslovaná, ale je uvedena v obsahu.
\def\chapwithtoc#1{
\chapter*{#1}
\addcontentsline{toc}{chapter}{#1}
}

\begin{document}

% Trochu volnější nastavení dělení slov, než je default.
\lefthyphenmin=2
\righthyphenmin=2

%%% Titulní strana práce

\pagestyle{empty}
\begin{center}

\large

Univerzita Karlova v Praze

\medskip

Matematicko-fyzikální fakulta

\vfill

{\bf\Large DIPLOMOVÁ PRÁCE}

\vfill

\centerline{\mbox{\includegraphics[width=60mm]{logo.eps}}}

\vfill
\vspace{5mm}

{\LARGE Bc. Ondřej Odcházel}

\vspace{15mm}

% Název práce přesně podle zadání
{\LARGE\bfseries Automatické doporučování ilustračních snímků}

\vfill

% Název katedry nebo ústavu, kde byla práce oficiálně zadána
% (dle Organizační struktury MFF UK)
Ústav formální a aplikované lingvistiky

\vfill

\begin{tabular}{rl}

Vedoucí diplomové práce: & RNDr. Pavel Pecina, Ph.D. \\
\noalign{\vspace{2mm}}
Studijní program: & Informatika \\
\noalign{\vspace{2mm}}
Studijní obor: & Matematická lingvistika \\
\end{tabular}

\vfill

% Zde doplňte rok
Praha 2014

\end{center}

\newpage

%%% Následuje vevázaný list -- kopie podepsaného "Zadání diplomové práce".
%%% Toto zadání NENÍ součástí elektronické verze práce, nescanovat.

%%% Na tomto místě mohou být napsána případná poděkování (vedoucímu práce,
%%% konzultantovi, tomu, kdo zapůjčil software, literaturu apod.)

\openright

\noindent
děkuji máje, že se nebojí létat

\newpage

%%% Strana s čestným prohlášením k diplomové práci

\vglue 0pt plus 1fill

\noindent
Prohlašuji, že jsem tuto diplomovou práci vypracoval(a) samostatně a výhradně
s~použitím citovaných pramenů, literatury a dalších odborných zdrojů.

\medskip\noindent
Beru na~vědomí, že se na moji práci vztahují práva a povinnosti vyplývající
ze zákona č. 121/2000 Sb., autorského zákona v~platném znění, zejména skutečnost,
že Univerzita Karlova v Praze má právo na~uzavření licenční smlouvy o~užití této
práce jako školního díla podle §60 odst. 1 autorského zákona.

\vspace{10mm}

\hbox{\hbox to 0.5\hsize{%
V ........ dne ............
\hss}\hbox to 0.5\hsize{%
Podpis autora
\hss}}

\vspace{20mm}
\newpage

%%% Povinná informační strana diplomové práce

\vbox to 0.5\vsize{
\setlength\parindent{0mm}
\setlength\parskip{5mm}

Název práce:
Automatické doporučování ilustračních snímků
% přesně dle zadání

Autor:
Bc. Ondřej Odcházel

Katedra:  % Případně Ústav:
Ústav formální a aplikované lingvistiky
% dle Organizační struktury MFF UK

Vedoucí diplomové práce:
RNDr. Pavel Pecina, Ph.D., Ústav formální a aplikované lingvistiky
% dle Organizační struktury MFF UK, případně plný název pracoviště mimo MFF UK

Abstrakt:
% abstrakt v rozsahu 80-200 slov; nejedná se však o opis zadání diplomové práce

Klíčová slova:
% 3 až 5 klíčových slov
vyhledávání obrazových informací

\vss}\nobreak\vbox to 0.49\vsize{
\setlength\parindent{0mm}
\setlength\parskip{5mm}

Title:
% přesný překlad názvu práce v angličtině
Automatic suggestion of illustrative images

Author:
Bc. Ondřej Odcházel

Department:
Institute of Formal and Applied Linguistics
% dle Organizační struktury MFF UK v angličtině

Supervisor:
RNDr. Pavel Pecina, Ph.D., Institute of Formal and Applied Linguistics
% dle Organizační struktury MFF UK, případně plný název pracoviště
% mimo MFF UK v angličtině

Abstract:
% abstrakt v rozsahu 80-200 slov v angličtině; nejedná se však o překlad
% zadání diplomové práce

Keywords:
% 3 až 5 klíčových slov v angličtině
information retrieval, image retrieval

\vss}

\newpage

%%% Strana s automaticky generovaným obsahem diplomové práce. U matematických
%%% prací je přípustné, aby seznam tabulek a zkratek, existují-li, byl umístěn
%%% na začátku práce, místo na jejím konci.

\openright
\pagestyle{plain}
\setcounter{page}{1}
\tableofcontents

%%% Jednotlivé kapitoly práce jsou pro přehlednost uloženy v samostatných souborech
\chapter{Úvod}
\addcontentsline{toc}{chapter}{Úvod}

Cílem diplomové práce je implementovat kompletní webovou aplikaci pro doporučování a vyhledávání ilustračních obrázků v textu. Vytořit takovou aplikaci přináší mnoho rozličných úkolů a problémů. Tato kapitola se bude snažit tyto problémy načrtnout. Další kapitola se bude jednotlivými problémy zabývat podrobně.

\section{Práce s daty}

Zadaná data obsahují 20 milionů anotací obrázků. Základním úkolem je být schopen takové množství dat vůbec nahrát do databáze a být schopný obsloužit mnoho požadavků za minutu. Bude zmíněn současný stav vývoje databázového software pro práci s velkými daty zejména s ohledem na snadnost hledání a škálovatelnost.

\section{Extrakce klíčových slov}

Extrakce klíčových slov je důležitý podobor NLP. V práci budou rozebrány algoritmy pro extrakci klíčových slov. Bude kladen zejména důraz na rychlost a nenáročnost na zdroje. Z uživatelských testovaní společnosti Google vychází, že rychlost načtení stránky je jedním z klíčových vlastností pro spokojenost uživatele. Klíčová slova budou mít v aplikaci dvě využití. Pokud uživatel zadá pouze text článku, extrahovaná klíčová slova se použijí na vyhledávání relevantních obrázků. Prvních několik klíčových slov bude navíc použita jako nápověda uživateli, ten pak může tato klíčová slova využít k exaktnímu omezení množiny klíčových slov.

\section{Překlad do češtiny}

Popisky klíčových slov jsou v angličtině. Tato práce řeší překlad množiny klíčových slov do češtiny. Kromě překladu je pro hledání také nutno implementovat algoritmus na stemming. Celá aplikace je navržena tak, aby případný další jazyk mohl být přidán co nejjednodušeji.

\section{Detekce jazyka}

Jednou z drobností, kterou ocení uživatel aplikace je detekce jazyka. Uživatel bude mít možnost zadat jazyk vstupního článku exaktně, ale aplikace bude také jazyk vstupního textu sama detekovat. Budou prozkoumány možnosti detekce jazyka. Opět se nejedná o nějakou klíčovou funkci aplikace. Výstup detekce bude moci být uživatelem měněn (podobně jako funguje Google Translate\footnote{https://translate.google.com/}), důraz bude tedy kladen na rychlost a jednoduchost.

\section{Webová aplikace}

Všechny předchozí komponenty se spojí v jedné webové aplikaci. Webový vývoj zažívá bouřlivý rozvoj. Na backendu jsou nové zejména způsoby práce s velkým množstvím dat v distribuovaném prostředí. Ve frontendové části probíhá rozvoj pomocí implementace nových technologií, známých pod hlavičkou HTML5, do moderních prohlížečů. Práce bude rozebírat všechny možnosti tvorby moderních webových aplikací.

\section{Testování}

Aplikace bude otestována na několika úrovních. Extrakce klíčových slov bude otestována pomocí korpusu článků a klíčových slov. Bude vytvořena komplexní webová aplikace pro testování doporučených obrázků. Tato aplikace bude vydělena ze samotné webové aplikace a bude používána i nezávisle.

\chapter{Zadání}

Většina zpravodajských serverů často opatřuje publikované články tzv. ilustračními snímky, jejichž úkolem je vizuálně dokreslovat obsah článku a upoutat na něj čtenářovu pozornost. Ilustrační snímky většinou pocházejí z rozsáhlých fotografických databází, jsou vybírány autory článku a s obsahem článku souvisejí jen relativně volně. Výběr ilustračních snímků probíhá nejčastěji na základě porovnávání klíčových slov specifikovaných autorem textu a popisků, kterými jsou obrázky v databázi opatřeny (typicky svými autory). 

Proces výběru ilustračních snímků (dotazování ve fotografické databázi) je obtížný jednak pro samotný vyhledávací systém (hledání relevantních fotografií na základě uživatelských dotazů), jednak pro autory, kteří musí dotazy vytvářet. Konstrukce dotazů spočívá v několika krocích: uživatel nejdříve musí identifikovat ústřední téma (či témata) článku, které chce ilustrovat vhodnou fotografií, a ta potom popsat vhodnými klíčovými slovy, zvolit a zkombinovat je tak, aby vedla k nalezení vhodného obrázku. Tento proces by mohl být zjednodušen tím, že konstrukce dotazů pro vyhledávání bude prováděna automaticky pouze na základě textu článku. 

Cílem diplomové práce je navržení a implementace komfortní webové aplikace pro automatické navrhování ilustračních snímků na základě textu článku, bez nutnosti explicitně konstruovat vyhledávací dotazy. Součástí práce bude i uživatelská evaluace celého systému. Pro experimenty bude použita kolekce ilustračních snímků od společnosti Profimedia.


\chapter{Název první kapitoly}

\section{Název první podkapitoly v první kapitole}

\section{Název druhé podkapitoly v první kapitole}


\chapter{Praktická část: Implementace moderní webové aplikace}

Práce je z velké části implementační. Snažil jsem se tedy vytvořit moderní webovou aplikaci s využitím co nejvíce frontendových novinek v novém standardu HTML5.

\section{Databáze: jak uložit 20M metadat obrázků}

Jak uložit takové množství dat, aby se dalo rychle vyhledávat. Škálovatelnost. Dostupnost knihoven pro práci s databázema. Proč jsem si vybral nakonec ES. SQL vs NoSQL.

\section{Backend a úprava dat: Komunikace s databází, implementace algoritmů}

Proč jsem zkoušel golang a proč jsem nakonec použil Ruby a Ruby on Rails. Základní popis MVC frameworku. Dostupnost knihoven pro získání stemů a prací s databází.

\section{Frontend: AJAXová aplikace na zobrazování obrázků}

Jaké jsou dnešní možnosti vývoje frontendu. Single page aplikace. Možnosti moderních prohlížečů. JavaScriptové knihovny. Proč to nedělám v jQuery, ale používám Google Closure. Google Closure Library, Templates, Compiler.

Návrh rozhraní bez jediného tlačítka. Responzivní webdesign.

\section{Anotační rozhraní}

Jak lze v Ruby on Rails vyrobit jednoduše anotační rozhraní s uživateli a s ukládáním do databáze.

\section{Návod k použití}

Popis prvků. Screenshoty aplikace.

\section{NOTES}

Language data jsem stahnul z:
http://dumps.wikimedia.org/cswiki/latest/

Ceska data jsou z:
/net/seznamdata/profiset/profi-text-cleaned.csv

prekladova data:
moses - http://www.statmt.org/moses/RELEASE-2.1/binaries/macosx-mavericks/bin/

prekladovy model:
http://www.statmt.org/moses/RELEASE-2.1/models/en-cs/


\chapter{Evaluace výsledků}

\section{Metodika}

Jak budu měřit. Porovnání naivního a nejlepšího algoritmu. Nejlepší algoritmus vybrán pomocí oanotovaných klíčových slov ve Wikinews článcích. Pak anotátoři se budou snažit najít nejméně vhodný obrázek z nabízených hodnot.

Evaluace jenom pro angličtinu, nebo i pro češtinu?

\section{Výsledky}

Grafy, tabulky.

\section{Možná zlepšení}

Jak bychom kmohli mít lepší data a co omezuje použitý algoritmus.

% Ukázka použití některých konstrukcí LateXu (odkomentujte, chcete-li)
%%%% Ukázka použití některých konstrukcí LaTeXu

\subsection{Ukázka \LaTeX{}u}
\label{ssec:ukazka}

V~této krátké části ukážeme použití několika základních konstrukcí \LaTeX{}u,
které by se vám mohly při psaní práce hodit.

Třeba odrážky:

\begin{itemize}
\item Logo Matfyzu vidíme na obrázku~\ref{fig:mff}.
\item Tato subsekce má číslo~\ref{ssec:ukazka}.
\item Odkaz na literaturu~\cite{lamport94}.
\end{itemize}

Druhy pomlček:
červeno-černý (krátká),
strana 16--22 (střední),
$45-44$ (minus),
a~toto je --- jak se asi dalo čekat --- vložená věta ohraničená dlouhými pomlčkami.
(Všimněte si, že jsme za~\verb|a| napsali vlnovku místo mezery: to aby se
tam nemohl rozdělit řádek.)

% Makro na české uvozovky (novější verze LaTeXu ho už mají zabudované)
\newcommand{\uv}[1]{\quotedblbase #1\textquotedblleft}
\uv{České uvozovky.}

\newtheorem{theorem}{Věta}
\newtheorem*{define}{Definice}	% Definice nečíslujeme, proto "*"

\begin{define}
{\sl Strom} je souvislý graf bez kružnic.
\end{define}

\begin{theorem}
Tato věta neplatí.
\end{theorem}

\begin{proof}
Neplatné věty nemají důkaz.
\end{proof}

\begin{figure}
	\centering
	\includegraphics[width=30mm]{../img/logo.eps}
	\caption{Logo MFF UK}
	\label{fig:mff}
\end{figure}


\chapter{Závěr}
\addcontentsline{toc}{chapter}{Závěr}

Podařilo se nám implementovat webovou aplikaci pro~vyhledávání ilustračních obrázků z~textu. Obrázky pochází z~datasetu Profimedie. Každý obrázek má přiřazena klíčová slova, která aplikace používá k~vyhledávání. Podporuje český a anglický jazyk, ale je snadno rozšiřitelná na podporu jiných jazyků. Uživatel může nejen zadat text článku, ale může si i vynutit přímo klíčová slova u hledaných obrázků.

Klíčovým algoritmem pro~vyhledávání je algoritmus extrakce klíčových slov. Ten je založen na modifikované metodě TF-IDF. Extrahovaná klíčová slova se používají k~vyhledávání doporučených obrázků. Druhé využití vyextrahovaných klíčových slov z~textu je jako nápověda uživateli.

Zdrojová data v~datasetu Profimedia jsou v~angličtině. Aby aplikace byla schopná pracovat s českými texty, bylo nutné klíčová slova obrázků přeložit. Použili jsme hybridní metodu, která kromě slovníkového překladu využívá i překlad detekovaných víceslovných frází v~textu. Menší kvalita takového překladu způsobuje občasné problémy algoritmu při doporučování obrázků k~českým textům. Při tak velkém množství dat by ovšem kvalitnější překlad byl velmi náročný~na zdroje, nebo příliš drahý.

Součástí práce je i rozbor některých moderních možností tvorby frontendu a backendu webových aplikací. Práce vysvětluje některé novinky v~HTML5 a ukazuje využití javascriptových frameworků pro~různé účely. Na backendové straně aplikace ukazuje některé nové komunikační a datové protokoly. Rozebrány jsou některé možnosti uložení velkého množství prohledavatelných dat.

Kromě textového vyhledávání umožňuje aplikace i vyhledávání podle vizuální podobnosti obrázků. Uživatel aplikace vidí v~detailu každého z~nalezených obrázků množinu podobných obrázků. v~rámci aplikace byla vytvořena nezávislá služba, která k~obrázkům v~datasetu Profimedia hledá obrázky. Podobnost obrázků je určena vzdáleností vektorů se $4\ 096$ složkami. Jelikož obrázků v~datasetu Profimedia je více než 20 milionů, bylo největším úkolem implementovat službu doporučování obrázků tak, aby vracela kvalitní výsledky co nejrychleji.

Algoritmus doporučování ilustračních obrázků byl testován s lidskými anotátory. Anotátoři dostávali české texty z~online médií a množinu obrázků. Některé obrázky byly získané algoritmem na doporučení obrázků k~textu. Jednu anotaci dostali nezávisle na sobě vždy dva anotátoři. v~prvním úkolu dostali anotátoři jeden náhodný a čtyři doporučené obrázky. z~58 anotací označili anotátoři správně náhodně vybraný výsledek ve~45 případech. Ukázalo se ovšem, že metoda testování má některé problémy, které umožnily anotátorům vybrat náhodný obrázek aniž by ostatní obrázky byly k~textu relevantní. Druhá metoda testování otočila poměr a anotace obsahovaly čtyři náhodné a jeden doporučený obrázek. Úkolem anotátorů bylo označit právě jeden nenáhodný obrázek. Každou anotaci opět dostali nezávisle na sobě vždy dva anotátoři. Ze 120 anotací se oba anotátoři shodli na správném obrázku v~76 případech. Rozbor anotací u kterých se oba anotátoři neshodli ukazuje různé příčiny. Některá klíčová slova k~obrázkům jsou chybně přeložena. Některé obrázky obsahují špatná klíčová slova. Klíčová slova některých obrázků jsou příliš krátká. v~několika anotacích by některé z~náhodně vybraných obrázků mohly být ilustrační obrázky k~textu také, což negativně ovlivnilo výsledky testování v~neprospěch algoritmu. pro~anglické články testování neproběhlo, ale je pravděpodobné, že dosažené výsledky by byly lepší --- u anglických klíčových slov nebylo potřeba použít automatický překlad.

Testování ukázalo, že aplikace ve~vysokém procentu nabízí uživatelům přijatelné ilustrační obrázky. Pokud jsou výsledky nevyhovující, může uživatel využít nápovědu doporučených klíčových slov a hledaný výsledek snadno přesněji specifikovat. ve~výsledku by tedy aplikace měla být schopna usnadnit proces vyhledávání ilustračních obrázků ke zpravodajským článkům, což bylo hlavním cílem práce.



%%% Seznam použité literatury
%%% Seznam použité literatury je zpracován podle platných standardů. Povinnou citační
%%% normou pro diplomovou práci je ISO 690. Jména časopisů lze uvádět zkráceně, ale jen
%%% v kodifikované podobě. Všechny použité zdroje a prameny musí být řádně citovány.

\def\bibname{Seznam použité literatury}
\begin{thebibliography}{99}
\addcontentsline{toc}{chapter}{\bibname}

\bibitem{lamport94}
  {\sc Lamport,} Leslie.
  \emph{\LaTeX: A Document Preparation System}.
  2. vydání.
  Massachusetts: Addison Wesley, 1994.
  ISBN 0-201-52983-1.

\bibitem{porter80}
  {\sc Porter,} Martin F.
  \emph{An algorithm for suffix stripping}.
  Program: electronic library and information systems, 1980, 14.3: 130-137



\end{thebibliography}


%%% Tabulky v diplomové práci, existují-li.
\chapwithtoc{Seznam tabulek}

%%% Použité zkratky v diplomové práci, existují-li, včetně jejich vysvětlení.
\chapwithtoc{Seznam použitých zkratek}

%%% Přílohy k diplomové práci, existují-li (různé dodatky jako výpisy programů,
%%% diagramy apod.). Každá příloha musí být alespoň jednou odkazována z vlastního
%%% textu práce. Přílohy se číslují.
\chapwithtoc{Přílohy}

\openright
\end{document}
